% Section 1: Introduction (compressed, 669→~550 words)

\section{Introduction}
\label{sec:introduction}

Large language models (LLMs) are increasingly used for information extraction (IE)~\cite{li2023codeie,sainz2024gollie,wei2023zeroshot}, but when generating relation triples $(h, r, t)$, neither entity necessarily carries a type consistent with the ontology's domain/range constraints for $r$.
We call such outputs \emph{domain/range violations} (DR): structurally valid triples that are ontologically invalid.
Violations are pervasive (15--18\% of triples under zero-shot decoding and 0.5--2\% even after LoRA fine-tuning) and correlate with schema restrictiveness rather than model scale.
At competitive quality, even low residual violation rates compound across thousands of documents into systematic knowledge graph corruption, making reliable constraint enforcement necessary.

Existing approaches leave a fundamental gap.
Schema-informed prompting reduces violations by 62\% on SciERC but is per-relation unreliable (violations for \textsc{Evaluate-for} actually \emph{increase} by 5.8\,pp).
Grammar-constrained decoding~\cite{dong2025xgrammar,willard2023outlines,beurer-kellner2023lmql} ensures syntactic well-formedness but cannot encode cross-field semantic constraints.
Post-hoc filtering discards up to 74\% of predictions including true positives (Appendices~\ref{app:cross-model-zs},~\ref{app:posthoc-extended}).
Training-based approaches~\cite{sainz2024gollie,lu2022uie} internalize schema knowledge but provide only soft guidance without runtime enforcement.

We propose \textbf{OntoCoDec}, a framework that translates IE ontologies into composable layers of decoding-time constraints.
A decidability analysis of IE type constraints (\S\ref{sec:decidability}) yields a two-layer architecture: the \emph{grammar layer} compiles ontology constraints into context-free grammars (CFGs) via XGrammar (JSON structure, vocabulary membership); the \emph{semantic layer} enforces type constraints at runtime via dynamic logit masking over an entity registry.
OntoCoDec compiles constraints automatically from IE ontologies, requires no constraint-specific training, and is applicable to any autoregressive LM with logits access.
We validate across three model families, two relation extraction (RE) datasets (up to 96 relation types), and event extraction (EE; \S\ref{sec:wikievents}), with 70B-class supporting evidence (\S\ref{sec:70b-results}).
Composed with LoRA, greedy decoding reduces DR to near-zero ($\leq$\,0.21\%) at relation F1 (R-F1) 40--46\% ($\pm$0.3 percentage points, $+$10\% latency), with 97.6--100\% of removed triples confirmed as false positives; beam ($k{=}4$) achieves near-zero violations at ${\sim}$1.5$\times$ latency.
DR enforcement converts irrecoverable ontology violations into cross-field type inconsistencies (CF;\,$\leq$\,1.4\% under fine-tuning), a severity downgrade amenable to post-hoc resolution.

OntoCoDec is orthogonal and complementary to training-based schema-aware approaches such as GoLLIE~\cite{sainz2024gollie} and KnowCoder~\cite{knowcoder_acl2024}; it can be composed with any generative IE system to improve schema compliance without retraining.

Our contributions:
\begin{enumerate}[leftmargin=*,itemsep=2pt]
    \item A two-layer constrained decoding framework that reduces domain/range violations to near-zero ($\leq$\,0.21\%) without retraining or significant quality loss, showing that these type constraints require stateful enforcement beyond context-free grammars (Proposition~1; \S\ref{sec:decidability}).
    \item An empirical finding of non-additive layer interaction: grammar-only reduces DR by 3\% and semantic-only by 24\%, yet their composition achieves 94\%, because grammar-enforced vocabulary membership creates the clean input the semantic layer requires.
    \item Scale-dependent layer interaction: the semantic layer's independent DR reduction grows from 24\% (8B) to 84\% (70B; \S\ref{sec:70b-results}), because larger models better satisfy the vocabulary prerequisites that enable semantic enforcement. Beam $k{=}4$ achieves near-zero violations at 8B; event extraction on WikiEvents confirms cross-task generalization (\S\ref{sec:wikievents}).
\end{enumerate}
